% Figure: schematic of a Ranque-Hilsch vortex tube. Compressed air enters
% tangentially, splits into a cold stream out one end and a hot stream out the
% other. No moving parts. The schematic shows the tangential inlet, the tube,
% and the two temperature-sorted outlets, with an entropy note that the sorting
% is paid for by the pressure drop of the supply air.
\begin{figure}[htbp]
\centering
\begin{tikzpicture}[>=Latex, font=\small]
% the tube body
\draw[thick, fill=enggray!8] (0,0) rectangle (8,1.6);
% cold outlet at left
\draw[thick, engblue, fill=engblue!10] (-1.4,0.3) rectangle (0,1.3);
\node[engblue, font=\scriptsize, align=center] at (-0.7,0.8) {cold\\end};
% hot outlet at right (a throttle cone)
\draw[thick, engred, fill=engred!10] (8,0) -- (9.2,0.2) -- (9.2,1.4) -- (8,1.6) -- cycle;
\node[engred, font=\scriptsize, align=center] at (8.75,0.8) {hot\\end};
% tangential supply inlet from below, near the cold end
\draw[thick, engamber, fill=engamber!12] (1.5,-1.2) rectangle (2.6,0);
\node[engamber, font=\scriptsize, align=center] at (2.05,-0.6) {supply\\(high $p$)};
% swirl arrows inside the tube
\foreach \x in {2.2,3.2,4.2,5.2,6.2} {
  \draw[enggray, ->] (\x,0.8) ++(-0.35,0) arc[start angle=180, end angle=-90, radius=0.35];
}
% cold stream arrow leaving left
\draw[engblue, very thick, ->] (-0.2,0.8) -- (-1.6,0.8);
\node[engblue, font=\scriptsize, anchor=south] at (-1.0,0.9) {$T_{\text{cold}}$};
% hot stream arrow leaving right
\draw[engred, very thick, ->] (9.2,0.8) -- (10.6,0.8);
\node[engred, font=\scriptsize, anchor=south] at (10.0,0.9) {$T_{\text{hot}}$};
% supply arrow entering
\draw[engamber, very thick, ->] (2.05,-1.5) -- (2.05,-1.25);
% caption note inside
\node[enggray, font=\scriptsize] at (4,2.1)
  {one supply stream splits into a cold stream and a hot stream, no moving parts};
\end{tikzpicture}
\caption[Schematic of a Ranque-Hilsch vortex tube]{The Ranque-Hilsch vortex
tube sorts a single stream of compressed air into a cold stream and a hot
stream with no moving parts. Air enters tangentially near one end at high
pressure and spins down the tube; the outer layers of the swirl leave hot
through a throttled cone at the far end, while the inner core reverses and
leaves cold through an orifice at the near end. The device looks, at first
glance, like a Maxwell demon made of plumbing, sorting fast molecules from
slow. It is not: the sorting is paid for by the large pressure drop of the
supply air, whose expansion generates far more entropy than the temperature
separation removes, so the cycle satisfies the second law with a wide margin.
The exergy analysis of the accompanying vignette shows the vortex tube to be a
thermodynamically poor cooler bought for its mechanical simplicity, the same
trade the throttle-versus-expander comparison drew for the household
refrigerator.}
\label{fig:vol04ch04:vortex-tube}
\end{figure}
